%% Self-consistent stationary noise floor for ZScoreES on quadratic landscapes.
%% Companion to theory.tex (Qiu et al. 2025).
%% Validated numerically in scripts/exp_E_ou_variance_trajectory.py.

\documentclass[american]{article}
\usepackage[T1]{fontenc}
\usepackage[utf8]{inputenc}
\usepackage{geometry}
\geometry{verbose,tmargin=2cm,bmargin=2cm,lmargin=2.5cm,rmargin=2.5cm}
\usepackage{amsmath}
\usepackage{amssymb}
\usepackage{amsthm}
\usepackage{hyperref}
\usepackage{babel}

\newtheorem{proposition}{Proposition}
\newtheorem{corollary}{Corollary}
\newtheorem{remark}{Remark}
\newtheorem{lemma}{Lemma}

\title{Self-Consistent Stationary Noise Floor for ZScoreES on Quadratic Landscapes}
\author{}
\date{2026-03-23}

\begin{document}
\maketitle

\section*{Setup and Notation}

We consider the ZScoreES update on the quadratic landscape
$R(\theta) = -\frac{1}{2}\theta^\top Q\theta$,
with $Q = U\Lambda U^\top$, $\Lambda = \mathrm{diag}(\lambda_1,\ldots,\lambda_d)$,
$\lambda_k \geq 0$.
Let $\mathbf{u}_k$ denote the $k$-th eigenvector of $Q$, and define the
eigenbasis coordinates $c_k := \mathbf{u}_k^\top\theta$.

From Proposition 3 of \texttt{theory.tex}, the per-step ES update satisfies
\begin{align}
    \mathbb{E}[\Delta\theta] &= \frac{-\alpha\sigma\, Q\theta}{\sigma_R(\theta)},
    \label{eq:mean-update}\\
    \mathrm{Cov}[\Delta\theta] &= \frac{\alpha^2}{N\sigma_R^2(\theta)}
        \Bigl(\sigma^2 Q\theta\theta^\top Q + 2\sigma^4 Q^2 + \sigma_R^2(\theta)\,I_d\Bigr),
    \label{eq:cov-update}
\end{align}
where $\alpha = \sigma/2$ and
\begin{equation}
    \sigma_R^2(\theta)
    := \sigma^2\|Q\theta\|^2 + \tfrac{\sigma^4}{2}\operatorname{Tr}[Q^2] + \sigma_\xi^2.
    \label{eq:sigmaR-def}
\end{equation}

The simplified linearized dynamics (frozen-$\sigma_R$ proposition in \texttt{theory.tex})
approximates $\sigma_R(\theta) \approx \sigma_R^0$ as a constant, giving the
discrete Ornstein--Uhlenbeck (OU) process
\begin{equation}
    \theta_{t+1} = \Bigl(I - \gamma^0 Q\Bigr)\theta_t + \eta_t,
    \qquad \eta_t \sim \mathcal{N}\!\Bigl(0,\,\tfrac{\alpha^2}{N}I_d\Bigr),
    \label{eq:frozen-ou}
\end{equation}
where $\gamma^0 = \alpha\sigma/\sigma_R^0$ is the frozen effective step size.
The stationary variance of each projected coordinate is
\begin{equation}
    V_k(\sigma_R^0)
    := \lim_{t\to\infty}\mathrm{Var}[c_k(t)]
    = \frac{\alpha^2/N}{1 - \gamma_k^{0\,2}},
    \qquad
    \gamma_k^0 = 1 - \frac{\alpha\sigma}{\sigma_R^0}\lambda_k,
    \label{eq:stationary-var-frozen}
\end{equation}
provided $|\gamma_k^0| < 1$.

The \emph{frozen-at-$\theta_0$} approximation sets $\sigma_R^0 = \sigma_R(\theta_0)$.
This overestimates the true stationary variance by a factor of $3$--$4\times$
when $\theta_0$ is far from $\theta^* = 0$, because $\sigma_R(\theta_0) \gg \sigma_R(\theta_\infty)$.

\section*{Main Result}

\begin{proposition}[Self-Consistent Stationary Noise Floor]
\label{prop:self-consistent}
Under the frozen-OU model \eqref{eq:frozen-ou}, the stationary distribution
satisfies $\theta_\infty \sim \mathcal{N}(0, \Sigma_\infty)$
with $\Sigma_\infty = \mathrm{diag}(V_1,\ldots,V_d)$ in the eigenbasis of $Q$.
The exact stationary noise floor is determined by the fixed-point equation
\begin{equation}
    \boxed{
    \sigma_R^{*\,2}
    = \sigma^2 \sum_{k=1}^d \lambda_k^2\, V_k(\sigma_R^*)
    + \tfrac{\sigma^4}{2}\operatorname{Tr}[Q^2]
    + \sigma_\xi^2,
    }
    \label{eq:self-consistent}
\end{equation}
where $V_k(\sigma_R^*) = \dfrac{\alpha^2/N}{1-(1-\frac{\alpha\sigma}{\sigma_R^*}\lambda_k)^2}$,
and $\sigma_R^*$ is the value of $\sigma_R$ consistent with the stationary distribution
it implies via \eqref{eq:sigmaR-def}.
\end{proposition}

\begin{proof}
At stationarity, $\theta_\infty \sim \mathcal{N}(0, \Sigma_\infty)$ with
$(\Sigma_\infty)_{kk} = V_k(\sigma_R^*)$ (from the OU stationary variance formula
\eqref{eq:stationary-var-frozen}, evaluated at $\sigma_R^*$).
The expected squared gradient norm at stationarity is
\begin{align}
    \mathbb{E}_{\theta \sim \mathcal{N}(0,\Sigma_\infty)}\bigl[\|Q\theta\|^2\bigr]
    &= \mathbb{E}\bigl[\theta^\top Q^2\theta\bigr]
    = \operatorname{Tr}[Q^2\,\Sigma_\infty]
    = \sum_{k=1}^d \lambda_k^2\, V_k(\sigma_R^*).
    \label{eq:E-Qtheta-squared}
\end{align}
Self-consistency requires that the value of $\sigma_R$ implied by the stationary
distribution \eqref{eq:sigmaR-def} equals the value $\sigma_R^*$ used to compute
the stationary variance:
\begin{align}
    \sigma_R^{*\,2}
    &= \sigma^2\,\mathbb{E}[\|Q\theta_\infty\|^2]
     + \tfrac{\sigma^4}{2}\operatorname{Tr}[Q^2]
     + \sigma_\xi^2 \notag \\
    &= \sigma^2 \sum_{k=1}^d \lambda_k^2\, V_k(\sigma_R^*)
     + \tfrac{\sigma^4}{2}\operatorname{Tr}[Q^2]
     + \sigma_\xi^2.
\end{align}
Substituting the expression for $V_k$ completes the proof.
\end{proof}

\begin{remark}[Connection to the frozen-at-$\theta_0$ approximation]
Setting $\sigma_R^0 = \sigma_R(\theta_0)$ rather than the self-consistent $\sigma_R^*$
inflates the denominator $1-\gamma_k^{0\,2}$ of \eqref{eq:stationary-var-frozen}
(because $\sigma_R(\theta_0) \gg \sigma_R^*$ when $\theta_0$ is far from $\theta^*$),
leading to a systematic overestimate of all $V_k$ and hence of the total noise floor
$\sum_k V_k$. Numerically, at $d=200$, $k=10$, $N=50$, $\sigma=0.1$:
$\sigma_R(\theta_0) = 0.758$, $\sigma_R^* = 0.044$, giving a $17\times$ change in
$\sigma_R$ and a $14\times$ overestimate of $\sum_k V_k$ by the frozen-at-$\theta_0$ formula.
\end{remark}

\begin{remark}[Existence and uniqueness of the fixed point]
Define the map $F: \mathbb{R}_{>0} \to \mathbb{R}_{>0}$ by
\begin{equation}
    F(s) := \sqrt{
        \sigma^2 \sum_{k:\,\lambda_k>0} \lambda_k^2\, V_k(s)
        + \tfrac{\sigma^4}{2}\operatorname{Tr}[Q^2]
        + \sigma_\xi^2
    }.
\end{equation}
For the stable regime ($0 < \lambda_k < 2\sigma_R^*/\alpha\sigma$ for all $k$),
$V_k(s)$ is a smooth, strictly decreasing function of $s > 0$
(larger $s$ $\Rightarrow$ smaller $\gamma_k$ $\Rightarrow$ larger $1-\gamma_k^2$
$\Rightarrow$ smaller $V_k$). Hence $F$ is strictly decreasing in the
signal-dominated region and satisfies $F(s) \to \sqrt{\tfrac{\sigma^4}{2}\operatorname{Tr}[Q^2]+\sigma_\xi^2}$
as $s \to \infty$ and $F(s) \to \infty$ as $s \to 0^+$.
By the intermediate value theorem, at least one fixed point $\sigma_R^* = F(\sigma_R^*)$ exists.
Uniqueness follows when $|F'(s)| < 1$ in a neighborhood of the fixed point,
which holds for moderate $\sigma$ and $N$ (verified numerically across all tested conditions).
\end{remark}

\begin{remark}[Practical computation]
The fixed point can be found by the damped iteration
\begin{equation}
    s_{n+1} = (1-\rho)\,s_n + \rho\,F(s_n),
    \qquad \rho \in (0,1),
\end{equation}
initialized at $s_0 = \sqrt{\tfrac{\sigma^4}{2}\operatorname{Tr}[Q^2] + \sigma_\xi^2}$
(the value at $\theta^* = 0$).
In practice, $\rho = 0.2$ and $\approx 60$ iterations suffice for convergence
to relative tolerance $10^{-8}$.
\end{remark}

\begin{corollary}[Self-consistent stationary variance per mode]
\label{cor:stat-var}
At the fixed point $\sigma_R^*$, the stationary variance of the $k$-th eigenmode is
\begin{equation}
    V_k^* := V_k(\sigma_R^*)
    = \frac{\alpha^2/N}{2\,\frac{\alpha\sigma}{\sigma_R^*}\lambda_k
            - \Bigl(\frac{\alpha\sigma}{\sigma_R^*}\Bigr)^2\!\lambda_k^2},
\end{equation}
and the total stationary noise floor is
\begin{equation}
    \mathbb{E}[\|\theta_\infty\|^2]
    = \sum_{k=1}^d V_k^*
    = \frac{\alpha^2}{N} \sum_{k=1}^d
        \frac{1}{2\gamma^*_k\lambda_k/(\sigma_R^*/\alpha\sigma)
              - (\gamma^*_k\lambda_k)^2/(\sigma_R^*/\alpha\sigma)^2},
\end{equation}
where $\gamma^*_k = 1 - (\alpha\sigma/\sigma_R^*)\lambda_k$.
In the well-conditioned limit $\gamma^*_k \approx 0$ (optimal curvature),
each mode contributes approximately $\alpha^2/N$, giving
$\mathbb{E}[\|\theta_\infty\|^2] \approx \alpha^2 k_{\mathrm{eff}}/N$
where $k_{\mathrm{eff}}$ is the number of stable active modes.
\end{corollary}

\section*{Numerical Validation}

The self-consistent equation \eqref{eq:self-consistent} was validated against
empirical simulations in \texttt{scripts/exp\_E\_ou\_variance\_trajectory.py}.

\paragraph{Setup.} Quadratic landscape with $d=200$, $k=10$ active modes,
power-law spectrum $\lambda_i = \lambda_{\max}\cdot i^{-1}$ ($\lambda_{\max}=5$,
$\lambda_{\min}=0.1$), $N=50$, $\sigma=0.1$, $T=300$ steps, $n_{\mathrm{trials}}=100$.

\paragraph{Results.}

\begin{center}
\begin{tabular}{lrr}
\hline
Quantity & Value & Notes \\
\hline
$\sigma_R(\theta_0)$ (frozen at $\theta_0$) & $0.75751$ & Initial value \\
$\sigma_R^*$ (self-consistent) & $0.04445$ & Fixed-point solution \\
$\sigma_R^{\mathrm{emp}}$ (empirical, $t=T$) & $0.04472$ & Measured from simulation \\
Relative error $|\sigma_R^* - \sigma_R^{\mathrm{emp}}|/\sigma_R^{\mathrm{emp}}$ & $0.6\%$ & \\
\hline
$\sum_k V_k^*$ (self-consistent) & $0.0027$ & \\
$\sum_k V_k^{\mathrm{emp}}$ (empirical) & $0.0030$ & \\
$\sum_k V_k^0$ (frozen at $\theta_0$) & $0.042$ & $14\times$ overestimate \\
\hline
\end{tabular}
\end{center}

The self-consistent formula matches the empirical noise floor to within $<1\%$,
compared to the $14\times$ overestimate from the frozen-at-$\theta_0$ approximation.

\paragraph{$\sigma$ dependence.}
The overestimate factor from the frozen-at-$\theta_0$ approximation grows as $\sigma$ decreases
(larger $\sigma$ $\Rightarrow$ $\sigma_R(\theta_0)$ is more dominated by the curvature term
$\tfrac{\sigma^4}{2}\operatorname{Tr}[Q^2]$, which does not change during optimization):

\begin{center}
\begin{tabular}{rrrrr}
\hline
$\sigma$ & $\sigma_R(\theta_0)$ & $\sigma_R$ final & $\sigma_R$ ratio & Frozen/emp \\
\hline
$0.05$ & $0.370$ & $0.011$ & $0.030$ & $27.8\times$ \\
$0.10$ & $0.758$ & $0.043$ & $0.057$ & $14.5\times$ \\
$0.50$ & $5.885$ & $1.118$ & $0.190$ & $4.7\times$ \\
\hline
\end{tabular}
\end{center}

\section*{Implications for theory.tex}

The existing proposition in \texttt{theory.tex} (Dynamics of ES on Quadratic Landscape)
correctly derives the frozen-OU variance trajectory and stationary variance for a
\emph{fixed} $\sigma_R$. The self-consistent equation \eqref{eq:self-consistent}
extends this to determine \emph{which} $\sigma_R$ is the physically relevant one:
it is not $\sigma_R(\theta_0)$ (appropriate near the start of optimization) but
$\sigma_R^*$ (the value consistent with the stationary distribution).

The correction adds no new parameters and requires only solving a single-variable
fixed-point equation in $\sigma_R^*$.

\end{document}
